\section{Conclusion}
\label{sec:conclusion}

Autonomous reasoning in uninstructed interactive environments is fundamentally constrained by perception rather than downstream algorithmic planning. Under the non-linear action penalties of ARC-AGI-3, agents cannot afford to compensate for impoverished representations through exploratory trial and error. As our empirical investigations demonstrate, presenting environments as unparsed coordinate matrices induces perceptual misinduction, compelling language and vision backbones to squander their action budgets attempting to reconstruct topological adjacency and basic spatial mechanics.

To resolve this limitation, we introduced \textsc{Sarbloh}, an autonomous agent architecture organized around the Measured--Guessed Rule. By computing deterministic geometric invariants directly on the host in linear time, the perception engine extracts lattice reachability, topological enclosures, midpoint tethers, and dihedral symmetries as mathematical certainties that generalize unconditionally across novel environments. Simultaneously, the framework treats latent operational semantics not as brittle classifications, but as provisional inductive hypotheses paired with minimal unit-action verification probes. Grounded by this perceptual foundation, the agent models environment transitions through programmatic Python execution, dispatching multi-step actuation trajectories that halt dynamically upon encountering unexpected state transitions.

Our findings indicate that advancing fluid intelligence on novel benchmarks does not strictly necessitate larger parameter scales or unconstrained inference budgets. Instead, equipping reasoning models with structured, core-knowledge perceptual abstractions bridges the divide between visual observation and symbolic deduction, enabling open-weight architectures to explore with disciplined, sample-efficient physical agency.
